\section{Comparison with the scalar local curve}\label{sec:triplecomparison}

Let $D\subset X$ satisfy Proposition~\ref{prop:isolation}, with $h=g(D)$
equal to three or four. Write $N=N_{D/X}$, and put
\[
                 H=3h+1,\qquad r_3=8,\qquad r_4=288 .
\]
Theorem~\ref{thm:classification} gives $r_h$ families of embedded
genus-$H$ curves in the total space of $N$. Write $\Phi$ for the map from complex to smooth families of
threefolds in the universal cycle algebra. We compare the scalar
insertion at $3[D]$ in this bundle with the compact open class $J_3$
of the isolated cycle in $X$. Write $G_k(X,3D)$ for the local virtual
degree of connected unmarked stable maps of arithmetic genus $k$ with
image cycle $3[D]$. It includes maps with contracted subcurves.
The inverse image of the isolated cycle is open and closed
in the degree-three stable-map space over a sufficiently small tube.
It is proper by fixed-genus compactness. The ordinary algebraic
stable-map virtual class restricts to it, using a projective completion
and the open set of maps with image in $X$.

\begin{proposition}\label{prop:borel}
Let $w=(2\sin(u/2))^2$, and let $S^{\rm sc}_{h,3}(u)$ be the connected
scalar residue in degree three. In the smooth cycle algebra, and in the
fixed-target connected Gromov--Witten theory, respectively, one has
\begin{align}
 \Phi(J_3)&=\Phi(x_{h,3})+\frac{r_h}{3}\Phi(x_{H,1}),
           \label{eq:cubic-smooth-comparison}\\
 \sum_k G_k(X,3D)u^{2k-2}
     &=S^{\rm sc}_{h,3}(u)+\frac{r_h}{3}w^{H-1}.
           \label{eq:cubic-residue-comparison}
\end{align}
\end{proposition}

\begin{proof}
\emph{The scalar insertion.}
Choose $p\in D$ outside the branch divisors of the $r_h$ families and
a linear form $\ell$ on $N_p$ nonzero on each of their defining
sublines. For a degree-three cycle let $F$ be the product of the
three values of $\ell$ over $p$, with multiplicities. On a family
$y^3=a s_i$ of Theorem~\ref{thm:classification}, this product is
$c_i a$ with $c_i\ne0$. Thus its unique zero is $3[D]$. Under
scalar dilation $F$ has weight three.

Take $B=\mathbb P^1$ and $L=\mathcal O_B(1)$. In the diagonal Borel
quotient a function of weight one is a section of $L$. Indeed a
pair $(v,z)$ defines the functional on $\mathbb Cv$ taking $v$ to
$z$, which is unchanged under simultaneous scaling. The function
$F$ therefore gives a section of $L^3$ on the relative cycle space.
The pullback of its Thom class divided by three, over the oriented
base $B$, represents $x_{h,3}$ by the scalar construction
\cite[\S4.3]{Pardon}.

We identify its algebraic virtual pairing before deforming the
insertion. Compactify the normal bundle by
$\mathbb P_D(\mathcal O\oplus N)$ and extend the supported class by
zero. The only fixed maps meeting its support are maps to $D$. On
their proper connected moduli stack
$\overline{\mathcal M}_k(D,3)$, with universal curve $\pi$ and map
$f$, the moving virtual normal complex is $R\pi_*f^*N$ with scalar
character one. If $\eta$ is the equivariant parameter, the
normalized test restricts to $3\eta/3=\eta$. The virtual localization
formula \cite{GP} gives the equivariant pushforward
\[
              \eta\,[u^{2k-2}]S^{\rm sc}_{h,3}(u).
\]
Here the inverse moving Euler class is exactly the fixed-target
residue of \cite[\S2.2]{BP}. It applies to $N$ since its
$K$-class on $D$ is determined by rank and determinant and its
scalar character is common to both normal directions. Passing to the
associated bundle over $B$ substitutes $c_1(L)$ for $\eta$, and
$\int_Bc_1(L)=1$. This computes the pairing of the relative
algebraic virtual class with the test in source genus $k$: on a
trivialization of the frame torsor the universal source, universal
map and their relative obstruction complex are pulled back from the
fixed-target equivariant data, and these identifications descend
with the scalar character.

\emph{The deformation.}
Form the deformation to the normal cone of $D\subset X$ and delete
the proper transform of $X\times\{0\}$. Denote its parameter by $t$.
The fibers over $0$ and over $t\ne0$ are $N$ and $X$, respectively.
Give $t$ weight $-1$ and fix the original $X$ coordinate. In normal
coordinates $n=y/t$ the action at $t=0$ is scalar of weight one.
The associated family of threefolds thus lies over
$E=\operatorname{Tot}_B(L^{-1})$. Its restriction to the zero
section is the preceding Borel family. Over the complement of that
section the original $X$ coordinate trivializes the family.

The test extends over a neighborhood of the marked cycle in this
family. Choose adapted local coordinates $(x,y_1,y_2)$ near $p$ with
$D=\{y_1=y_2=0\}$ and take the transverse surface $x=x(p)$.
For a nearby cycle its three intersections with this surface near $p$,
counted with multiplicity, give a continuous unordered triple. In
the deformation coordinates use the product of their values under
$\ell(y/t)$. The divided coordinates extend over $t=0$, and the
product has weight three. It restricts there to $F$ and vanishes on
the marked triple in every fiber. On smooth embedded curves transverse
to the surface the evaluation is holomorphic.

Choose a small section $v\in H^0(B,L^3)$ with three simple zeros.
Choose a small smooth section $\tau$ of $L^{-1}$ with a single
transverse zero $b_*$ disjoint from those of $v$. The metric dual of
a holomorphic section of $L$ with a simple zero gives such a section
of local index $-1$. On the zero section first deform $F$ to $F-v$.
Keep this translated test fixed and move the parameter over the
graphs $t=s\tau(b)$, $0\le s\le1$.

These intersections have compact support through both homotopies.
Use a closed normal disk of small radius over $D$. At $t=0$ the
maximum normal radius on the $i$th curve is a positive constant
times $|a|^{1/3}$. Hence $|F|$ is bounded below away from zero
where a cycle reaches the boundary of this disk, uniformly over
$B$ in unitary frames and over the finitely many families. Choose
$\|v\|$ smaller than this bound. If $t\ne0$, the disk maps into
an $O(|t|)$ tube about $D$ in the original target. For $|t|$
sufficiently small its only degree-three cycle is $3[D]$ by
Proposition~\ref{prop:isolation}. Within the closed disk, a common
closed taming form in each local smooth trivialization bounds the
masses by the fixed homology class $3[D]$. Relative cycle compactness
on a slightly larger disk therefore includes any limit across
$t=0$ \cite[Theorem~2.11]{DIW}. These choices can be made uniform
on the compact base. In particular the compactly supported
cohomology class on cycles is transported in the smooth colimit
\cite[\S2]{WueProduct}. For the algebraic pairing use the relative
virtual class over the smooth algebraic base $E$ and intersect it
with the Thom class of the oriented parameter graph and the test.
The inverse images of the compact supports are proper at each fixed
$k$. The homotopic graphs give equal pairings. At the zero section
the refined fiber restriction is the virtual class of the Borel
family by base change \cite[Proposition~7.2]{BF}.

\emph{Intersections on the final graph.}
Over each of the three zeros of $v$, the parameter $\tau$ is nonzero.
The only supported cycle is the isolated triple in $X$. Its local
intersection is the simple zero of $-v$, of positive index. It gives
$\Phi(J_3)$ in the smooth colimit and $G_k(X,3D)$ in the relative
virtual pairing, since the family is a product there.

Over $b_*$ there is one solution $a=v(b_*)/c_i\ne0$ on each
normal family. The local relative Hilbert scheme is smooth with
coordinates $(b,a)$ and maps to $(b,0)\in E$. Indeed the only
fixed-fiber tangent direction is the variation of $a$, and the
$t$ direction surjects onto the one-dimensional fixed-fiber
obstruction space by Proposition~\ref{prop:isolation}. The
normal curves themselves supply the smooth two-dimensional family
in $t=0$. On this chart the graph equation is $-\tau(b)=0$,
and the remaining equation has nonzero complex derivative $c_i$
in $a$. Their intersection sign is $-1$. The graph has the
orientation of the vertical complex line in $E$, and the test
has its complex orientation. Interchanging their normal factors
preserves orientation.

The Hilbert chart here satisfies the hypotheses of
Proposition~\ref{prop:moving-embedded-curve}. In particular the
motion of $t$ surjects onto the fixed-curve obstruction and the
graph is transverse before testing the amplitude. The shear argument
of Lemma~\ref{lem:signed-embedded-cycle} is local along the embedded
image. At the finitely many supporting points choose disjoint smooth
submersion coordinates away from the evaluation and ramification
fibers. These coordinates extend over a smaller parameter neighborhood
by smooth local trivialization. The required normal variations are
compactly supported there; a projection from the whole threefold to
the original curve is unnecessary. On the cycle space the lemma gives
$-\Phi(x_{H,1})$ for each such image. The proper degree-one map stack
over the same chart accounts for the virtual contribution in every
source genus, equal by \eqref{eq:embedded-degree-one} to
$-[u^{2k-2}]w^{H-1}$. The embedded image has trivial automorphisms
as a map to the threefold; contracted components of maps to it retain
their ordinary stack multiplicities.

Summing the two sets of local intersections before the scalar
normalization gives
\[
        3\Phi(J_3)-r_h\Phi(x_{H,1}),\qquad
        3G_k(X,3D)-r_h[u^{2k-2}]w^{H-1}.
\]
After division by three these are respectively $\Phi(x_{h,3})$
and the coefficient of $u^{2k-2}$ in $S^{\rm sc}_{h,3}$ computed
above. This proves both identities.
\end{proof}
Let $\Phi$ denote the change from complex families to smooth families
of almost-complex threefolds in the Chern-number-zero universal
algebra. We recall the detection results used here
\cite[\S\S5--7]{WueProduct}. In the complex algebra a primitive class
annihilated by every proper division of cycles belongs to the finite
rational span of the $x_{j,1}$. Their images under $\Phi$ are linearly
independent. The latter assertion uses proper pullback to connected
stable maps of increasing arithmetic genus. In genus below $j$ that
pullback is zero for $x_{j,1}$, and in genus $j$ the regular embedding
gives a nonzero class. Neither assertion requires an evaluation by a
virtual fundamental class for arbitrary analytic families.

\begin{proof}[Proof of the universal identity in Theorem~\ref{thm:isolated}]
Let $J_0=x_{h,0}=1$. On the tube provided by
Proposition~\ref{prop:isolation}, the inverse image of the supported
point $d[D]$, $1\le d\le3$, under cycle addition is the union of the
pairs $(a[D],b[D])$ for $a+b=d$. The preimage under multiplication
by $m$ is the divided point if $m\mid d$, and is empty otherwise.
The maps pull back unit cochains, so there are no stabilizer factors.
Thus the $J_d$ and the scalar $x_{h,d}$ have the same coproduct and
division laws through degree three. In particular
\[
 \Delta J_d=\sum_{a+b=d}J_a\otimes J_b,\qquad
 \rho_m J_d=
 \begin{cases}
 J_{d/m}&m\mid d,\\
 0&m\nmid d .
 \end{cases}
\]
For the scalar classes these formulas follow from
\cite[\S4]{Pardon} in Chern number zero.

For $d=1,2$ deformation to the normal bundle preserves the isolated
section $d[D]$. Indeed a sequence of additional cycles would, by normal
dilation and compactness as in \S\ref{sec:neighborhood}, give an
off-section cycle of degree at most two in the linear normal model.
The entire degree-$d$ cycle space of that model is a point, and the
empty tuple represents $x_{h,d}$. Consequently
$\Phi(J_d)=\Phi(x_{h,d})$. At the first degree of a possible difference
the displayed operations make $J_d-x_{h,d}$ primitive and killed by
proper divisions. Detection therefore gives $J_d=x_{h,d}$ successively
in degrees one and two. For $d=3$ subtract also $(r_h/3)x_{H,1}$.
This correction is primitive and killed by all proper divisions.
The difference lies in the same multiplicity-one span, and has zero
smooth image by Proposition~\ref{prop:borel}. This proves
\eqref{eq:main-class}. The numerical identity is the separate
algebraic pairing in that proposition.
\end{proof}


In the compact example, change just the cubic coefficient of the
genus-three scalar factor to $J_3$.
Its logarithm gains $(8/3)x_{10,1}$. The new exponent at $3\alpha$
in genus ten is therefore $c_{3\alpha,10}-(4720/3+8/3)c_{\alpha,3}
=c_{3\alpha,10}-1576c_{\alpha,3}$. Thus the isolated local unit
accounts for the fractional part even though the complete change
from equal to opposite weights also has an integral part.

In genus five the quotient by dilation contains pencils. A single
cubic evaluation no longer suffices; two independent point products
give the least tuple in Theorem~\ref{thm:cubic-evaluations}. Theorem~\ref{thm:pencil}
in Appendix~\ref{app:pencils} proves the comparison in this case,
including the virtual multiplicity of a cuspidal normalization.
